\documentclass{article}
\usepackage{graphicx} % Required for inserting images
\usepackage[utf8]{inputenc} % Per caratteri accentati
\usepackage[T1]{fontenc}    % Per una corretta codifica dei font
\usepackage[italian]{babel} % Per sillabazione in italiano
\usepackage{float}
\usepackage{hyperref}
\usepackage{listings}
\usepackage{color}
\usepackage{amssymb}
\usepackage{amsmath}
\usepackage{multicol}

%VALORI SISTEMABILI
\usepackage[a4paper, top= 1 cm, bottom= 1.5 cm, left= 2 cm, right= 2 cm]{geometry}

\title{Sistemi Operativi}
\author{Francesco Lanzarini}
\date{ }

\begin{document}

\maketitle

\section{Domande}
\subsection{Ristorante}
\subsubsection{Cosa è un sistema operativo?}
Un sistema operativo e' un programma, piu' precisamente e' il progrmamma che viene lanciato all'accensione e terminato allo spegnimento.
\subsubsection{Cosa è un programma?}
Un programma e' una sequenza di istruzioni.
\subsubsection{Cosa è un processo?}
Un processo e' un programma che viene eseguito.
\subsubsection{Come è strutturato il modello di macchina di Von Neumann?}
Il modello di macchina di Von Neumann e' composto da processore che esegue le operazioni, una memoria che fornisce gli operandi delle operazioni, 
una memoria di massa dove memorizzare informazioni a lungo termine, bus dati per il trasferimento delle informazioni e dispositivi di esterni che gestiscono l'input/output.
\subsubsection{Cosa è un processore?}
Il processore e' il componente del calore che esegue le operazioni di fetch-decode-execute.
\subsubsection{Cosa è un bus?}
Un bus e' un canale di informazioni condiviso collegato tra piu' componenti, il suo compito e' quello di tyrasportare i dati.
\subsubsection{Cosa è un controller di Input/Output?}
Un controller I/O è un componente che gestisce lo scambio di dati tra un sistema di controllo e i dispositivi esterni.
\subsubsection{Cosa è un device register?}
Un device register è la registrazione di un dispositivo in una rete o in un sistema, associandolo a un identificativo univoco per riconoscerlo e gestirlo
\subsubsection{Cosa è un device driver?}
Un device driver è un software che permette al sistema operativo di comunicare con un dispositivo hardware, traducendo le sue richieste in comandi comprensibili dal dispositivo
\subsubsection{Cosa può fare un processore quando è in kernel mode?}
Richiedere/fornire informazioni  ai bus dati
\subsubsection{Cosa può fare un processore quando è in user mode?}
Eseguire le istruzioni dei programmi che non siano SO
\subsubsection{Come e in che casi avviene il passaggio da user mode in kernel mode?}
Il passaggio avviene principalmente in tre casi: system call, interrupt opppure trap/exceptions
\\subsubsection{Cosa è una system call?}
La system call e' l'interfaccia con cui un programma in User Mode chiede al sistema operativo un servizio kernel.
\subsubsection{Cosa è un Process Control Block?}
Un PCB e' una struttura dati che serve al sistema operativo per la memorizzazione delle informazioni dei processi.
\subsubsection{Cosa vuole dire scheduling?}
Per scheduling si intende l'operazione del sistema operativo atta a gestire l'esecuzione da parte della CPU di processi "ready".
\subsubsection{Come vengono sospesi i processi in attesa di operazioni di I/O?}
I processi vengono sospesi grazie ad una system call.
\subsubsection{Cosa significa multitasking?}
Per multitasking si intende la capacita' della CP di eseguire piu' processi contemporaneamente.
\subsubsection{Cosa è un interrupt e quando avviene?}
Un interrupt (interruzione) è un segnale inviato al processore per interrompere il flusso di esecuzione corrente
e richiedere l'esecuzione immediata di una speciale routine del kernel, puo' avvenire per diversi motivi tipo la digitazione di un carattere.
\subsubsection{Una system call o un interrupt fa passare il controllo dal processo in esecuzione al sistema operativo. 
Cosa deve fare il processore per consentire di poter riprendere successivamente l'esecuzione del processo dal punto dell'interruzione?}
Per consentire la corretta proseguita del processo il sistema operativo ha a disposizione tutti i dati necessari nella PCB del singolo processo.
\end{document}